% !TeX program = pdflatex
\documentclass[11pt,a4paper]{article}

% --- Język i czcionki ---
\usepackage[T1]{fontenc}
\usepackage[utf8]{inputenc}
\usepackage[polish]{babel}
\usepackage{lmodern}
\usepackage[final]{microtype}

% --- Układ strony i nagłówki/stopki ---
\usepackage[a4paper,margin=2.5cm]{geometry}
\usepackage{fancyhdr}
\pagestyle{fancy}
\fancyhf{} % czyść wszystko
\lhead{Notatki z matematyki}
\rhead{\leftmark}
\cfoot{\thepage}

% --- Matematyka ---
\usepackage{amsmath,amssymb,amsthm,mathtools}
\usepackage{siunitx} % jeżeli potrzebne jednostki
\sisetup{locale = PL} % separator 1,23 (polski), można zmienić na locale=PL gdy dostępne
\usepackage{bm}       % pogrubione symbole

% --- Odsyłacze i tytuły ---
\usepackage[hidelinks]{hyperref}
\usepackage[nameinlink,capitalise,noabbrev]{cleveref}
\usepackage{csquotes}

% --- Listy i drobiazgi ---
\usepackage{enumitem}
\setlist{noitemsep,topsep=3pt}
\usepackage{xcolor}

% --- Środowiska twierdzeń po polsku ---
\newtheoremstyle{polski}% name
  {6pt}{6pt}% space above/below
  {\itshape}% body font
  {}% indent
  {\bfseries}% head font
  {.}% punctuation
  {0.5em}% spacing after head
  {}% head spec
\theoremstyle{polski}
\newtheorem{tw}{Twierdzenie}[section]
\newtheorem{lem}[tw]{Lemat}
\newtheorem{wniosek}[tw]{Wniosek}

\newtheoremstyle{polski-def}
  {6pt}{6pt}{\normalfont}{}{\bfseries}{.}{0.5em}{}
\theoremstyle{polski-def}
\newtheorem{defi}[tw]{Definicja}
\newtheorem{przyklad}[tw]{Przykład}
\newtheorem{uwaga}[tw]{Uwaga}

% --- Nazwy do cleveref po polsku ---
\crefname{tw}{twierdzenie}{twierdzenia}
\Crefname{tw}{Twierdzenie}{Twierdzenia}
\crefname{lem}{lemat}{lematy}
\Crefname{lem}{Lemat}{Lematy}
\crefname{wniosek}{wniosek}{wnioski}
\Crefname{wniosek}{Wniosek}{Wnioski}
\crefname{defi}{definicja}{definicje}
\Crefname{defi}{Definicja}{Definicje}
\crefname{przyklad}{przykład}{przykłady}
\Crefname{przyklad}{Przykład}{Przykłady}
\crefname{uwaga}{uwaga}{uwagi}
\Crefname{uwaga}{Uwaga}{Uwagi}
\crefname{equation}{równanie}{równania}
\Crefname{equation}{Równanie}{Równania}
\crefname{section}{sekcja}{sekcje}
\Crefname{section}{Sekcja}{Sekcje}

% --- Szybkie skróty ---
\newcommand{\N}{\mathbb{N}}
\newcommand{\Z}{\mathbb{Z}}
\newcommand{\Q}{\mathbb{Q}}
\newcommand{\R}{\mathbb{R}}
\newcommand{\C}{\mathbb{C}}
\DeclareMathOperator{\sgn}{sgn}
\DeclareMathOperator{\Var}{Var}
\DeclareMathOperator{\E}{\mathbb{E}}

% --- Numeracja i wygląd sekcji ---
\numberwithin{equation}{section}
\setcounter{secnumdepth}{3}
\setcounter{tocdepth}{2}

% --- Meta ---
\title{Sprawozdanie z etapu 1.}
\author{Jakub Kogut}
\date{\today}

\begin{document}
\maketitle
\section{Członkowie zespołu}
Nasz zespół składa się z 5 członków:
\begin{itemize}
    \item Jakub Kogut -- kierownik zespołu,
    \item Manfred Gawlas,
    \item Szymon Hładyszewski,
    \item Wiktor Koczkodaj,
    \item Wojciech Typer,
\end{itemize}

\section{Temat i zakres projektu}
\textbf{Temat:} Aplikacja mobilna do nauki o inwestowaniu, przeznaczona dla studentów Uniwersytetu Ekonomicznego we Wrocławiu. \newline
\textbf{Zakres:}
\begin{itemize}
    \item Aplikacja mobilna na system Android, zawierająca:
        \begin{itemize}
            \item moduł z materiałami edukacyjnymi: wykładami (statyczne)
            \item moduł z quizami sprawdzającymi wiedzę (statyczne/możliwe dynamiczne)
            \item forum do dyskusji dla użytkowniów
            \item system rankingu użytkowników (otrzymywane punkty za quizy)
        \end{itemize}
    \item API do zarządzania danymi aplikacji
        \begin{itemize}
            \item zarządzanie użytkownikami
            \item zarządzanie materiałami edukacyjnymi
            \item zarządzanie quizami
            \item zarządzanie forum
            \item zarządzanie rankingiem
            \item panel administratora
        \end{itemize}
        Zintegrowane z bazą danych
\end{itemize}

\section{Wykorzystane technologie}
Planujemy wykorzystać następujące technologie:
\begin{itemize}
    \item Aplikacja mobilna: Kotlin + Jetpack Compose
    \item API: python + Django + Django REST framework
    \item Baza danych: PostgreSQL (obsługiwane przez Django)
\end{itemize}

\section{Zarządzanie projektem}
Do zarządzania chcemy wykorzystać Github Projects. Repozytorium dostępne jest pod adresem: \url{https://github.com/niooch/bizwise.git}
            
\end{document}

